\section{La carrera de la IA, el código abierto y el panorama de la industria}

\subsection{¿Quién está ganando dinero?}

La única empresa que realmente está ganando dinero es NVIDIA. Las hiperescaladoras muestran utilidades contables, pero su gasto de capital es enorme. OpenAI va a la cabeza en el mejor modelo y en los ingresos de IA, pero gasta más de lo que ingresa. Meta se beneficia de la IA a través de sus sistemas de recomendación (no directamente de Llama).

\begin{warningbox}{La velocidad de la comoditización de las capacidades}
\textit{«Cualquier empresa cuyo modelo de negocio se basara en capacidades del nivel de GPT-3 está muerta. Cualquiera que se basara en el nivel de GPT-4 también está muerta.»} (Dylan)

OpenAI y Anthropic tienen que seguir ganando en el mejor modelo, o serán sustituidos por Llama o por modelos de código abierto. La curva de caída de costos es asombrosa: la inferencia de GPT-3 pasó de \$60-70 por millón de tokens a unos cuantos centavos de dólar en la actualidad, una caída de 1,200 veces en 3 años. DeepSeek está sobre la misma línea de tendencia de precios de los modelos del nivel de GPT-4; solo fue el primero en llegar.

Esto significa que el modelo de negocio de «modelo como servicio» enfrenta una presión de fondo. El único foso competitivo es mantenerse por delante de la frontera de manera continua: en cuanto te atrasas medio año, un modelo de código abierto alcanza tus capacidades y las vuelve gratuitas.
\end{warningbox}

\subsection{Agent: exageración vs. realidad}

«Agent» es el término de IA más abusado de 2025; su significado real debería ser la resolución de tareas abierta e independiente.

\begin{knowledgebox}{El problema de los seis nueves}
Cada paso tiene una precisión $<$100\%, y al componer varios pasos el rendimiento final cae mucho. Si cada paso tiene 99\% de precisión, después de 10 pasos solo queda 90\%; después de 100 pasos solo queda 37\%. ¿Cuántos nueves necesita cada paso para alcanzar una tasa de éxito de extremo a extremo de 99.99\%?

\textit{«¿Cuántos nueves tienes? Multiplícalos por el número de pasos... ni siquiera 99.9999\% alcanza.»} (Dylan)

La analogía con la conducción autónoma: incluso en carreteras bien definidas sigue siendo enormemente difícil (a Waymo le tomó diez años); un entorno abierto de páginas web o sistemas operativos es varias veces más caótico. Cada sitio web tiene un diseño distinto, las API cambian en cualquier momento y los patrones de error son innumerables.
\end{knowledgebox}

\begin{figure}[H]
\centering
\includegraphics[width=0.92\textwidth]{frame-05-agents.jpg}
\caption{Discusión sobre los límites de viabilidad y confiabilidad de los Agents\protect\footnotemark}
\end{figure}
\footnotetext{Intervalo de tiempo en el video: 04:17:00--04:17:12.}

La dirección más prometedora es el \textbf{Agent de ingeniería de software}, por las siguientes razones:
\begin{itemize}
    \item Verificabilidad: las pruebas unitarias, la compilación y el CI/CD ofrecen retroalimentación automática
    \item Puede revisar toda la base de código (la ventana de contexto es suficiente)
    \item SWE-bench pasó de 4\% a 60\% en un año
\end{itemize}

La colaboración estructurada también funciona (OpenAI Operator + DoorDash/OpenTable): en dominios bien definidos, la tasa de éxito de los Agent es mayor.

El costo de la ingeniería de software va a desplomarse, lo cual cambia toda la estructura del mercado. China no usa SaaS porque sus ingenieros son baratos, y la IA llevará ese modelo a todo el mundo. Los expertos de dominio (aviación, semiconductores, industria química) usan herramientas de hace 20 años: \textit{«Las herramientas de litografía de ASML corren en Windows XP.»} (Dylan) Estos ámbitos tienen mucha fruta al alcance de la mano.

\subsubsection{DOGE y la modernización del gobierno}

El software del gobierno está sumamente atrasado: \textit{«suplica una modernización»}. La burocracia protege centros de poder; el software rompe esas barreras. Muchas industrias tienen abundante fruta al alcance de la mano para la automatización con IA.

\subsection{La controversia de la destilación}

\begin{knowledgebox}{La zona gris ética de la destilación}
La destilación (entrenar un modelo más débil con las salidas de uno más fuerte) es una práctica estándar en la industria. OpenAI afirma que DeepSeek usó las salidas de su API, pero esos datos ya estaban en internet por medio de copiar y pegar. Muchas empresas emergentes estadounidenses entrenan abiertamente con salidas de OpenAI.

\textit{«¿Por qué es poco ético que yo entrene con las salidas de tu modelo, mientras que tú puedes entrenar con el texto de internet?»} (Nathan)

El espionaje industrial también existe: las ideas fluyen libremente a través de los cambios de empleo en Silicon Valley (en California las cláusulas de no competencia son ilegales). Un ingeniero de Gemini que ayudó a construir el context de 1 millón se fue después a Meta; se espera que la siguiente generación de Llama tenga un context de 1 millón.

El resquicio de Japón: la ley de derechos de autor permite entrenar con cualquier dato, más 9 GW de energía nuclear detenida, más importación ilimitada de GPU, es igual a un posible paraíso para el entrenamiento de IA.
\end{knowledgebox}

\subsection{El impacto de DeepSeek R1 en el código abierto}

La licencia MIT de R1 es un «gran reinicio»: el primer modelo de frontera con una licencia verdaderamente permisiva. Antes, las opciones eran modelos que no eran de frontera o licencias restrictivas (como la de Llama).

\textit{«Necesitamos modelos verdaderamente abiertos... no sabemos cuáles son los datos de DeepSeek R1, pero puedes usarlo para hacer una copia barata y hacer como si fuera tuya. Con Llama no puedes hacer eso.»} (Nathan)

El proyecto Tulu de AI2 de Nathan demuestra el poder de un posentrenamiento completamente abierto: al aplicar RLVR (RL con recompensas verificables) sobre la base de Llama, superó a Llama instruct en matemáticas e igualó a DeepSeek V3. Entre los modelos ubicados entre los primeros 60 lugares de Chatbot Arena, antes de Tulu \textbf{ninguno} había publicado el código o los datos de su posentrenamiento.

\subsection{Resumen de la sección}

El panorama de la industria de la IA se está consolidando con rapidez: solo NVIDIA es rentable, y la comoditización de las capacidades avanza a toda velocidad. Los Agent enfrentan el reto de confiabilidad de los seis nueves. El código abierto está cambiando las reglas del juego: la licencia MIT de DeepSeek R1 y el proyecto Tulu de AI2 demuestran que un posentrenamiento completamente abierto puede alcanzar el nivel de frontera. La controversia de la destilación refleja un problema sin resolver sobre la propiedad de los datos.
